% ---------- common preamble : CM3191 organic report drafts (XeLaTeX) ----------
\usepackage[a4paper,margin=22mm]{geometry}
\usepackage{fontspec}
\usepackage{xeCJK}
\setmainfont{TeX Gyre Termes}
\setsansfont{TeX Gyre Heros}
\setCJKmainfont{Noto Serif CJK KR}
\xeCJKsetup{CJKspace=true, CJKecglue={}}
\usepackage{amsmath,amssymb}
\usepackage[version=4]{mhchem}
\usepackage{chemfig}
\usepackage{siunitx}
\sisetup{detect-all, per-mode=symbol}
\usepackage{booktabs,tabularx,array}
\usepackage{enumitem}
\usepackage{xcolor}
\usepackage[super]{cite}
\usepackage{titlesec}
\usepackage[hidelinks]{hyperref}
\usepackage{caption}
\usepackage{float}
\captionsetup{font=small,labelfont=bf}
\titleformat{\section}{\large\bfseries}{\thesection.}{0.5em}{}
\titleformat{\subsection}{\normalsize\bfseries}{\thesubsection}{0.5em}{}
\setlength{\parskip}{3pt}
\setlength{\emergencystretch}{3em}
\sloppy
\linespread{1.1}
\setchemfig{atom sep=1.9em, bond style={line width=0.6pt}}

% ---- placeholders: everything in a grey box must be replaced after the lab
\newcommand{\res}[1]{\fcolorbox{black}{gray!18}{\footnotesize\sffamily #1}}
\newcommand{\obs}[1]{\fcolorbox{black}{gray!18}{\footnotesize\sffamily OBS: #1}}
\newcommand{\resp}[1]{\par\smallskip\noindent\fcolorbox{black}{gray!18}{\parbox{\dimexpr\linewidth-2\fboxsep-2\fboxrule\relax}{\footnotesize\sffamily #1}}\par\smallskip}

% ---- title block: \reporttitle{title}{experiment line}{date}
\newcommand{\reporttitle}[3]{%
\begin{center}
{\LARGE\bfseries #1}\\[4mm]
{\large Kim Donghwan (A0266094X)}\\[1mm]
CM3191 Chemical Synthesis Experiments, #2\\
Thursday laboratory, fume hood B7 \quad\textbullet\quad Date of experiment: #3\\
Instructor: Dr Hoang Truong Giang
\end{center}
\vspace{2mm}
\noindent\rule{\linewidth}{0.6pt}}

% ---- Korean drafting notes page (removed with \draftfalse)
\newcommand{\draftnotes}[2]{%
\ifdraft
\thispagestyle{empty}
\begin{center}{\Large\bfseries 작성 안내 (제출 전 이 페이지 삭제)}\end{center}
\vspace{2mm}
\textbf{이 초안의 구조.} 결과와 무관한 부분(서론·이론, 실험 방법, 고찰의 이론적 틀, 참고문헌)은 완성된 문장으로 써 두었습니다. 실험 후 바뀌는 값과 관찰은 \res{회색 상자}로 남겼습니다. 상자 안의 숫자는 매뉴얼의 계획값이므로 \textbf{반드시 본인 측정값으로 바꿔야 합니다}. 결과에 따라 문장이 달라지는 곳은 .tex 파일에 \texttt{\%} 주석으로 선택지를 달아 두었습니다.

\textbf{실험 후 채울 것}
\begin{enumerate}[leftmargin=6mm,itemsep=0.5mm]
#1
\end{enumerate}

\textbf{제출 전 확인할 것 [CONFIRM]}
\begin{enumerate}[leftmargin=6mm,itemsep=0.5mm]
\item 유기 파트 보고서 형식·분량·마감(평가 방식 공지 확인). proforma 형식이면 문단을 답안으로 옮겨 쓰면 됩니다.
\item 과목의 생성형 AI 사용 규정(사용 사실 표기 필요 여부).
#2
\end{enumerate}
\textbf{컴파일}: XeLaTeX. 파일 첫 줄 근처의 \texttt{\textbackslash drafttrue}를 \texttt{\textbackslash draftfalse}로 바꾸면 이 페이지가 빠집니다.
\newpage
\setcounter{page}{1}
\fi}
